% collatz_spektrum_entwurf.tex
% Populärwissenschaftlicher Entwurf (Spektrum-Stil)
% Kompilierung: pdflatex collatz_spektrum_entwurf.tex

\documentclass[11pt,a4paper]{article}

\usepackage[utf8]{inputenc}
\usepackage[T1]{fontenc}
\usepackage[ngerman]{babel}
\usepackage{lmodern}
\usepackage{microtype}
\usepackage[a4paper,margin=2.5cm]{geometry}
\usepackage{parskip}
\usepackage{amsmath,amssymb}
\usepackage{booktabs}
\usepackage{xcolor}
\usepackage{titlesec}
\usepackage[most]{tcolorbox}
\usepackage[hidelinks]{hyperref}

\definecolor{spektrumblue}{RGB}{0,85,150}

\titleformat{\section}{\large\bfseries\color{spektrumblue}}{\thesection.}{0.5em}{}

\tcbset{
  infobox/.style={
    colback=blue!4,colframe=spektrumblue,fonttitle=\bfseries,
    arc=2pt,boxrule=0.8pt,left=6pt,right=6pt,top=4pt,bottom=4pt},
}

\newcommand{\Ztwo}{\mathbb{Z}_2}

\begin{document}

\begin{center}
  {\LARGE\bfseries Die Collatz-Vermutung: Zwischen Spielerei und tiefster Struktur}\\[0.6em]
  {\normalsize Entwurf für \emph{Spektrum der Wissenschaft}}\\[1em]
  {\normalsize Thomas Hoffbauer}\\[1.2em]
\end{center}

\noindent\textbf{Lead.}
Seit fast neunzig Jahren multiplizieren Mathematiker ungerade Zahlen mit drei,
addieren eins, halbieren so oft es geht --- und wissen nicht, warum jede Bahn
bei eins enden sollte. Ein strukturiertes Forschungsprogramm ordnet die Dynamik
nicht mehr dem Fixpunkt, sondern der Geometrie der 2-adischen Zahlen $\Ztwo$ zu.

\section*{Kurzfassung}

Die Collatz-Regel ist trivial; der Beweis ist es nicht. Terence Tao zeigte 2019,
dass \emph{fast alle} Startzahlen gutartig sind --- doch ``fast alle'' ist nicht
``alle''. Das Projekt kartiert die odd-to-odd-Dynamik über eine EABC-Klassifikation
modulo~12, beweisbare C-Ketten-Struktur und einen Paradigmenwechsel: Uniformität
als Kontraktion gegen eine Ausnahmemenge $E\subset\Ztwo$, nicht als Nähe zu~$1$.
\textbf{Die Vermutung bleibt offen.}

\begin{tcolorbox}[infobox,title={Was ist Lean?}]
Lean ist ein \emph{Beweisassistent}, kein Rechner, der Milliarden Collatz-Schritte
durchprobiert. Mathematische Aussagen werden darin wie Programme geschrieben;
jeder Beweisschritt wird maschinell geprüft oder als explizit offen markiert.
Lean beglaubigt damit, was wirklich bewiesen ist --- und was formal widerlegt ist.
Für die Collatz-Forschung trennt es zuverlässig Strukturtheorie von Vermutung.
\end{tcolorbox}

\begin{tcolorbox}[infobox,title={Was ist formalisiert?}]
\begin{center}
\renewcommand{\arraystretch}{1.25}
\begin{tabular}{@{}p{0.42\textwidth}p{0.48\textwidth}@{}}
\toprule
\textbf{Status} & \textbf{Inhalt} \\
\midrule
Bewiesen (Lean) &
mod-12-Struktur, C-Ketten-Endlichkeit, LTE-Kompensation,
2-adische Metrik $\mathrm{dist}_2$, Ausnahmemengen $E_N$ \\
Formal widerlegt &
Uniformität über $\mathrm{dist}_2(n,1)$; pauschale $\nu_2$-Schranken \\
Numerisch / heuristisch &
Blockdrift $\Lambda\approx -0{,}830$, mod-12-Mischung \\
\textbf{Offen} &
Limes $E\subset\Ztwo$; Uniformität $\mathrm{dist}_2(T^k(n),E)\to 0$;
volle Collatz-Vermutung \\
\bottomrule
\end{tabular}
\end{center}
\end{tcolorbox}

\section*{Historische Einordnung}

Lothar Collatz formulierte die Vermutung in den 1930er Jahren. Paul Erd\H{o}s
soll gesagt haben, die Mathematik sei noch nicht reif dafür. Tao (2019) bewies
asymptotische Gutartigkeit für fast alle Startwerte --- ein Durchbruch, der die
punktweise Lücke nicht schließt. Numerische Tests bis weit über $10^{20}$ finden
kein Gegenbeispiel; das ersetzt keinen Beweis.

\section*{Struktur statt Rätsel}

Die EABC-Klassifikation (Reste $1,5,7,11$ mod~12) verwandelt Bahnen in Wörter
über $\{\mathrm{E},\mathrm{A},\mathrm{B},\mathrm{C}\}$. Bewiesen sind unter anderem
Endlichkeit von C-Ketten, Verbot von B$\to$B und erzwungene EA-Schritte nach
maximalen C-Ketten. Falsch gewesene Hoffnungen --- globale Schranken über
$\mathrm{dist}_2(n,1)$ --- sind formal widerlegt und grenzen den Suchraum ein.

\section*{Offenes Ende}

Die Collatz-Vermutung bleibt ein Spiegel zwischen Spiel und Tiefe. Was das Projekt
liefert, ist eine präzise kartographierte Landschaft: formal abgesicherte
Strukturtheorie mit klar isolierter Uniformitätslücke an $E\subset\Ztwo$.
Ob dieser geometrische Ansatz den letzten Schritt von ``fast alle'' zu ``alle''
ermöglicht, ist ungewiss --- die Frage ist erstmals geometrisch präzise gestellt.

\vspace{0.8em}
\noindent\small\textit{Der Autor arbeitet an der formalen Strukturtheorie der
odd-to-odd Collatz-Dynamik. Eine ausführliche Darstellung liegt als Preprint vor;
die Collatz-Vermutung bleibt offen.}

\end{document}
